\section{Outcome Verification: partiendo de GSM8K}

\subsection{Training Verifiers to Solve Math Word Problems}

El trabajo de Cobbe et al. aporta al mismo tiempo el conjunto de datos GSM8K y un método sistemático para entrenar un verifier. GSM8K contiene alrededor de 8.5K problemas de matemáticas de primaria en forma de texto, cada uno acompañado de un proceso de razonamiento en lenguaje natural y una respuesta numérica final.

\videofigure{figures/gsm8k-contributions.jpg}{GSM8K y el outcome verifier constituyen las dos aportaciones centrales de este trabajo.}{00:02:10--00:04:26}

El proceso central es: primero se entrena el generator; para cada problema de entrenamiento se muestrean numerosas completions; se asigna una etiqueta binaria según si la respuesta final coincide con el ground truth; después se entrena el verifier para predecir si toda la solución es correcta.

\videofigure{figures/generator-verifier-pipeline.jpg}{El generator produce candidatos etiquetados y el verifier aprende a puntuar la solución completa.}{00:04:26--00:06:20}

\subsection{El objetivo conjunto del Verifier}

El verifier puede añadir una cabeza de clasificación sobre un modelo de lenguaje base, conservando al mismo tiempo el objetivo de next-token:

$$
\mathcal L(\phi)
=\mathcal L_{\mathrm{LM}}(\phi)
+\lambda\,\mathcal L_{\mathrm{cls}}(\phi),
$$

$$
\mathcal L_{\mathrm{cls}}
=-z\log v_{\phi}(x,y)-(1-z)\log(1-v_{\phi}(x,y)).
$$

\begin{itemize}
    \item $\phi$: parámetros del verifier;
    \item $z\in\{0,1\}$: si la respuesta final es correcta;
    \item $v_{\phi}(x,y)$: probabilidad de que el candidato sea correcto;
    \item $\lambda$: peso relativo entre el modelado de lenguaje y la clasificación de corrección.
\end{itemize}

\videofigure{figures/verifier-objective.jpg}{El verifier modela a la vez el texto y la corrección de la solución completa.}{00:06:20--00:08:16}

\begin{knowledgebox}{Por qué se conserva el objetivo de modelado de lenguaje}
Una clasificación pura podría aprender solo atajos como el formato, la longitud o palabras clave de la respuesta final. El objetivo de modelado de lenguaje obliga a la representación a seguir entendiendo el texto completo de razonamiento, y cuando los datos son limitados también actúa como regularización. Pero no elimina automáticamente los shortcuts; aun así es necesario revisarlos mediante ejemplos adversarios y evaluación fuera de distribución.
\end{knowledgebox}

\subsection{Los datos de entrenamiento deben provenir de la distribución de errores propia del modelo}

Si el verifier solo observa errores evidentes construidos manualmente, no necesariamente logrará reconocer los patrones de fallo reales del generator. Este método muestrea numerosas completions del modelo para cada problema, de modo que el conjunto de entrenamiento incluya negativos difíciles que «parecen mucho una solución correcta».

\begin{warningbox}{El Verifier no puede obtener una puntuación alta solo con datos estáticos}
Una mejora del generator, un cambio de prompt o un aumento en la longitud del razonamiento modifican la distribución de errores. Un verifier preciso sobre candidatos antiguos puede fallar con un modelo nuevo. Por ello, los datos del verifier deben muestrearse de manera continua a partir del generator actual, y evaluarse con control de versiones.
\end{warningbox}

\subsection{Best-of-N y el cómputo en tiempo de inferencia}

En el momento de la inferencia se generan $N$ candidatos y se elige el que obtiene la puntuación más alta del verifier:

$$
y^*=\arg\max_{y_j\sim\pi_\theta(\cdot\mid x)} v_\phi(x,y_j).
$$

A medida que $N$ aumenta, mejora la coverage de candidatos; mientras el ordenamiento del verifier conserve poder de discriminación, la precisión final también mejora. Se observa experimentalmente que el desempeño puede seguir mejorando hasta varios cientos de completions.

\videofigure{figures/testtime-compute.jpg}{Aumentar el número de candidatos sigue mejorando la precisión tras la selección del verifier.}{00:16:50--00:19:04}

\subsection{La limitación fundamental de outcome supervision}

La etiqueta de la respuesta final no puede indicar en qué paso se equivocó el razonamiento. Una trayectoria puede ser completamente incorrecta al principio y, por coincidencia, llegar a un valor numérico correcto; o bien el proceso puede ser en lo esencial correcto y solo fallar en la aritmética al final. El ORM etiqueta ambos casos, respectivamente, como positivo y negativo, con una asignación de crédito muy tosca.

\subsection{Resumen de la sección}

El outcome verifier demostró por primera vez que la learned verification puede convertir el inference sampling en una mejora real del desempeño, y constituye la base de la posterior RL reward. Pero atender solo a la respuesta final confunde el proceso correcto, el acierto casual y el error localizado, lo que impulsó a la investigación hacia la process supervision.
